\chapter{Introduction}
\label{ch:introduction}

This guide helps you prepare Git and GitHub for coursework. It covers setup, secure access, publishing your project, and submitting updates. Follow the chapters in order, replace examples such as \texttt{USERNAME} and \texttt{REPOSITORY}, and follow your instructor's repository requirements.

Git tracks changes on your computer; GitHub hosts an online copy of your project. A commit saves a snapshot of your work, and pushing sends it to GitHub for review.

\subsection*{The workflow}
\begin{center}
\begin{tikzpicture}[
  guideIcon/.style={circle,draw=UTPBlue!24,fill=UTPPaleBlue!38,line width=0.6pt,minimum size=1.05cm,inner sep=0pt},
  guideArrow/.style={-{Stealth[length=2mm]},draw=UTPTeal!65,line width=0.8pt}
]
  \draw[guideArrow] (0.55,0) -- (4.6,0);
  \draw[guideArrow] (5.7,0) -- (9.75,0);
  \node[guideIcon] (prepare) at (0,0) {{\large\color{UTPBlue}\faDownload}};
  \node[guideIcon] (secure) at (5.15,0) {{\large\color{UTPBlue}\faKey}};
  \node[guideIcon] (publish) at (10.3,0) {{\large\color{UTPBlue}\faGithub}};
  \node[align=center,text width=4.1cm] at (0,-1.4) {
    {\bfseries\color{UTPNavy}01\quad Prepare}\\[1mm]
    {\footnotesize Install Git and set your commit identity}
  };
  \node[align=center,text width=4.1cm] at (5.15,-1.4) {
    {\bfseries\color{UTPNavy}02\quad Connect securely}\\[1mm]
    {\footnotesize Add an SSH key and verify GitHub access}
  };
  \node[align=center,text width=4.1cm] at (10.3,-1.4) {
    {\bfseries\color{UTPNavy}03\quad Publish}\\[1mm]
    {\footnotesize Create a repository, push your project, and submit updates}
  };
\end{tikzpicture}
\end{center}

\subsection*{What you will learn}
\begin{itemize}[leftmargin=1.5em,itemsep=0.3em,topsep=0.2em,label=\textcolor{UTPTeal}{\faCheck}]
  \item Install Git and configure the name and email recorded in your commits.
  \item Create an SSH key, register its public half, and test your connection safely.
  \item Initialize a local repository, link it to GitHub, and push the first commit.
  \item Check, commit, and push later changes so your submitted work stays current.
\end{itemize}

\subsection*{Common mistakes to avoid}
\begin{itemize}[leftmargin=1.5em,itemsep=0.3em,topsep=0.2em,label=\textcolor{red!70!black}{\faTimes}]
  \item Never add passwords, tokens, API keys, \texttt{.env} files, or your private SSH key to a repository.
  \item Do not commit until \texttt{git status} shows only the files you intend to submit.
  \item Deleting a secret in a later commit does not remove it from earlier repository history.
\end{itemize}

\begin{tcolorbox}[
  enhanced,
  colback=UTPPaleBlue!18,
  colframe=UTPBlue!20,
  boxrule=0.5pt,
  arc=2mm,
  left=4mm,right=4mm,top=3mm,bottom=3mm
]
\noindent{\bfseries\color{UTPNavy}Before you begin.}
\begin{enumerate}[leftmargin=1.7em,itemsep=0.55em,topsep=0.55em,parsep=0.1em]
  \item \textbf{Check access.} Sign in to GitHub and confirm that you can access your account.
  \item \textbf{Install and configure Git.} Follow the installation and commit-identity instructions in Chapter~\ref{ch:git-ssh-setup}.
  \item \textbf{Confirm course requirements.} Check your instructor's repository naming and visibility requirements.
\end{enumerate}
\end{tcolorbox}